\section{서론}

% --------------------------------------------------------------------
% 이 파일의 역할
% Purpose of this file
% --------------------------------------------------------------------
% 1-Introduction.tex에는 논문의 서론을 작성합니다.
% 연구의 필요성, 연구 목적, 연구 문제, 용어 정의를 쓰는 곳입니다.
% Write the introduction of the thesis in this file.
% This chapter usually includes the research background, purpose,
% research questions, and definitions of key terms.

\subsection{연구의 필요성}

과학교육 분야의 학위논문에서 연구의 필요성은 독자에게 왜 이 과학 교수ㆍ학습 문제를 연구해야 하는지 설득하는 부분이다.
교육과정 변화, 과학 개념 학습의 어려움, 탐구 역량과 과학적 소양에 대한 요구, 학교 현장에서 관찰되는 수업 문제처럼 연구 주제가 등장한 교육적 맥락을 제시한다.
그다음 선행 연구가 무엇을 밝혀 왔고, 아직 충분히 설명되지 않은 지점이 무엇인지 정리한다.

\subsection{연구 목적}

연구 목적은 연구의 필요성에서 제기한 문제에 대해 이 논문이 무엇을 밝히려는지 선언하는 부분이다.
연구 대상, 과학 개념이나 탐구 기능, 교수ㆍ학습 처치, 학습자 특성, 분석 방향이 드러나도록 구체적으로 쓴다.

\subsection{연구 문제}

연구 문제는 연구 목적을 실제 자료로 답할 수 있는 질문으로 나눈 것이다.
보통 2--4개 정도로 제시하고, 각 연구 문제는 연구 방법의 분석 절차와 연구 결과 장의 하위 절에 대응되도록 구성한다.

\begin{enumerate}
    \item 학습자는 해당 과학 개념이나 탐구 기능에 대해 어떤 이해와 어려움을 보이는가?
    \item 제안한 교수ㆍ학습 처치는 학습자의 개념 이해나 탐구 수행에 어떤 변화를 가져오는가?
    \item 학습자 특성이나 수업 맥락에 따라 그 변화 양상은 어떻게 다르게 나타나는가?
\end{enumerate}

\subsection{용어 정의}

연구에서 반복해서 사용하는 핵심 개념이나 조작적 정의가 필요한 용어를 이 절에서 정리한다.
예를 들어 특정 이론에서 빌려 온 용어, 이 연구에서만 좁혀 쓰는 의미, 선행 연구마다 다르게 쓰이는 용어는 이 연구에서 사용하는 정의를 명시해 독자와의 오해를 줄인다.
